\documentclass[10pt]{article}
\usepackage{template/template}
\NewsSetKeys{date={2026年11月18日},vol={036},title={每日有趣新闻},author={编辑部},note={待验收稿：新闻核查与数学预印本检索已进行；尚未完成原始图片归档、PDF构建和视觉验收，不得作为正式发行版。}}
\begin{document}
\begin{documents}
\newsitem{AI热潮开始遭遇资本与能源双重约束}{深度报道｜经济、科技与气候｜11月17日}{人工智能基础设施投资仍在扩大，但市场越来越关心一个不同于模型能力的问题：建设数据中心所投入的资本，能否在可接受的期限内产生收入，电力系统又能否承受这一扩张？美国银行11月基金经理调查显示，受访者对科技股的净超配比例从上月37\%降至12\%，为2006年6月以来最大的单月降幅；38\%认为企业投资过度，52\%预期大型云计算公司半年内会宣布削减资本支出。调查覆盖169名管理约4610亿美元资产的投资者。这些数字反映的是样本内的投资者预期，而不是全市场投资者一致认定泡沫破裂，更不是企业已经实际削减了一半投资。\cite{bofa}

资本市场的疑虑与实体能源约束正在交汇。土耳其举行的COP31推出《安塔利亚人工智能承诺》，鼓励采用清洁电力、提高基础设施效率、披露排放并帮助发展中国家获得技术。据路透社11月17日报道，启动时约有90个签署者，但美国主要超大规模云服务商没有参加，主要科技企业中阿里巴巴参与。自愿倡议不是国际法义务，签署数量也不能代表排放量覆盖率；不参加亦不自动证明企业没有任何减排措施。报道引述的研究称美国400余座超大规模数据中心的二氧化碳排放与丹麦全国相近，这一比较还须注意核算边界、统计年度和电力间接排放的归属。\cite{cop}

为什么算力与电力不能简单换算？训练和推理需要GPU等计算设备，设备功耗只是用电的一部分。冷却、配电损耗、网络和储存都需要能源，通常以总设施耗电与IT设备耗电之比，即PUE，描述机房的能源开销。但PUE低并不意味着绝对排放低：一座规模增长数倍的高效数据中心仍可能增加总耗电。电网的碳强度又随地区、时段及发电组合变化。购入可再生能源证书、签订长期购电协议和每小时真实匹配低碳电源，在实际减排含义上并不等价。

对投资者而言，设备采购订单、已签合同、实际建成的电力容量、客户付费推理需求和自由现金流也属于不同阶段。高额订单既可能说明企业看好需求，也可能反映电力与设备交期紧张下的提前占位。市场调查预期企业削减开支，不意味着已建项目立刻消失；反过来，数据中心竣工也不能保证盈利。资金成本、芯片更新周期、设备折旧与能源合同期限，都决定了最终回报。

这也是为什么AI的宏观经济叙事需要分清短期与长期。建设期购买芯片、变压器、冷却设备与发电设施会直接刺激某些产业，但这种需求未必意味着所有行业劳动生产率同步跃升。美联储近期讨论了AI对就业、金融稳定及生产率的不同可能路径；这些属于情景研究，而非已证实的长期增长预测。编辑判断是，观察AI经济应当同时追踪算力利用率、真实收入、资本支出和用电排放，而不能只看模型跑分或科技股市值。\cite{bofa,cop,barr}}

\newsitem{NASA Terra卫星将结束科学观测：一颗设计寿命六年的卫星工作了二十七年}{重点报道｜航天与地球科学｜11月17日}{NASA科学新闻11月17日宣布，Terra地球观测卫星将在12月1日停止收集科学数据。它于1999年发射，最初设计任务寿命约六年，却持续工作二十七年。Terra属于地球观测系统的重要平台，长期同时观察大气、陆地、海洋与生物圈。对研究者而言，卫星最宝贵的资产不只是某一张漂亮的地球照片，而是连续多年采用相近仪器与处理流程获得的时间序列。\cite{terra}

长时间序列能让研究者区分季节波动、短期异常与长期变化，但仪器老化、轨道变化、定标方法升级也会引入观测偏差。卫星停止采集后，历史资料并不会因此失去价值；相反，跨卫星数据拼接和重新标定的重要性会提高。不能把卫星退役直接解释成地球观测能力全面中断，NASA还有其他在轨观测平台；也不能假定新仪器测得的数据无需校准即可与旧资料直接相连。

截至本期截稿，NASA所宣布的是未来停止科学数据采集的计划，而非11月17日已经停止运行。正式配图应选取Terra原始任务资料或经过标注的观测图，明确其观测日期、仪器与处理级别；由于尚未取得并归档可核查原始图片，本短刊不放未经确认许可的网页图片。}

\newsitem{数学拾趣：一个图上需要几名警察才能抓住强盗？}{数学与Lean｜11月17日新预印本}{图论中的Cops and Robbers是一类轮流行动的完全信息博弈。给定一张图，警察与强盗各占一个顶点，沿边移动或原地不动；警察的目标是最终与强盗占据同一顶点。所谓cop number，是保证抓捕成功所需的最少警察数。它不是顶点数或边数的简单函数，而与图的全局连接方式和双方策略有关。

Aarav Ghai于11月17日出现在arXiv数学新稿列表的论文研究flower snark与广义Petersen图，作者主张对所有$n\geq5$的flower snark $J_n$，cop number为3，而$J_3$为2；还声称完成了一项关于广义Petersen图何时具有cop number 2的猜想。上界通过作者称为rolling wall的三警察策略实现，下界则要构造强盗逃避两名警察的策略。文章摘要还称，相关定理及新的下界已经在Lean 4中形式化验证。\cite{ghai}

这里的验证至少有三层：首先，作者在预印本中声称了哪些定理；其次，Lean内核究竟检查了哪些形式命题及其证明；最后，形式命题是否忠实表达论文中的数学对象和博弈规则。Lean通过不等于我们已经独立审查了全部自然语言论证。本期只核对了arXiv列表上的作者、题名、摘要与提交状态，未取得形式化仓库并复现检查，也未确认经过同行评议。因此将其作为有趣的研究线索，而不是已由编辑部验证的定理公告。}

\newsitem{APEC讨论安全使用AI：共识与约束仍有距离}{国际外交与科技治理｜11月17日}{据路透社11月17日报道，亚太经合组织高级官员会议主席陈旭表示，各成员就安全、负责任地部署人工智能进行了深入讨论，并将数据对AI发展的支持作用纳入议题。报道发生在中美科技竞争加剧、美国进一步限制中国企业通过境外数据中心获得先进AI芯片的背景下。\cite{apec}

需要区别会议主席对讨论成果的概括、各成员实际通过的文件，以及具有强制执行力的法律规定。APEC通常以经济合作与政策协调为核心，成员对数据流通、产业安全和技术出口的利益并不一致。报道中的“重要共识”不能自动推导为统一的数据监管标准或对芯片出口限制的共同立场。相关政策的真正影响仍需追踪正式文件、各经济体国内执行办法和企业实际合规成本。}
\newsitem{编辑检索说明}{11月17日—18日}{本期已核查国际政治、经济、AI、气候、航天与arXiv数学新稿。其他领域未取得足够有价值且可交叉核实的当日材料，不以旧闻凑数。图像尚未下载归档，故本稿暂不使用任何未经确认许可的配图。}
\end{documents}
\small\section*{参考资料}\bibliographystyle{plain}\bibliography{sources/2026/Nov/18/refs}
\end{document}